\section{实验协议与验证重点}
\subsection{双臂任务的完成判据}
评价首先在参考视角定义完整目标区域足迹 $R^\star$，
并令 $\Omega$ 为该参考视角的整幅图像域。
令 $O_t^\star$ 为目标完整投影，$C_t$ 为遮挡布料投影。
当目标完整可评估且 $|O_t^\star|>0$ 时，定义
\begin{align}
 p_{\rm in}(t)&=\ind\{O_t^\star\subseteq R^\star\},\\
 a_C(t)&=|C_t\cap\Omega|/|\Omega|,\\
 p_{\rm restore}(t)&=\ind\{a_C(t)\ge\theta_C\}.
\end{align}
$R^\star$ 不是当前可见region类别mask：
互斥分割使object与可见region不重叠，直接求交会误判。
固定目标区域可由无遮挡参考帧独立标注；相机或区域移动时须重新估计。
仅有局部可见的目标mask也不能替代 $O_t^\star$ 来证明整个物体已进入区域。

当前布料复原\emph{不会重新遮挡目标物体}，复原指标是布料在画面中的面积占比，
不是目标被遮挡比例，也不是默认裁剪桌面ROI后的占比。
按这一要求，任务在规定时限内满足下式即可完成：
\begin{equation}
 \Phi_\ell(z_{0:T})=
 \ind\{\exists\,t\le T:
 p_{\rm in}(t)=1\ \land\ p_{\rm restore}(t)=1\}.
 \label{eq:completiontask}
\end{equation}
两项必须在同一完成时刻成立：不能仅凭初始布料面积较大，
也不能凭曾经入区但随后又移出的记录判成功。
若协议另要求连续 $d_C$ 帧确认，则在同一窗口内检验这两个谓词，
并要求窗口终点不超过 $T$；未指定持续时间时不默认添加这一要求。
当前 $T,\theta_C,d_C$ 的数值待统一评测协议确定，
不从历史其他实验的50\%、80\%或95\%阈值推定。

若最终物体与目标区域均可清晰评价，式\eqref{eq:completiontask}的两个谓词
可以由同一参考帧判定，不必依赖假设的最终目标遮挡。
执行中被机械臂临时遮挡仍可能导致评价证据不足，但不能把这种局部情况
混同于任务要求布料最后重新遮住物体。
暴露、分离、推动和复原是有助于解释行为的过程事件；
并非所有任务都必须采用同一组过程标签，更不要求在问题定义中指定阶段头。

\subsection{已有结果与需要解释的机制}
\label{sec:evidence}
研究首先由实机现象驱动。表\ref{tab:preliminary}记录用户目前提供的近似成功率，
而非本稿重新汇总的最终统计结果。
\begin{table}[t]
 \centering
 \caption{初步实机结果；次数、视角及配对条件待逐项核对}
 \label{tab:preliminary}
 \begin{tabular}{lr}
 \toprule
 方法描述 & 近似成功率\\
 \midrule
 基础ACT & 50\%\\
 ACT + MLP残差 & 60\%\\
 ACT + GRU残差 & 70\%\\
 语义输入策略（无QToken） & 70\%\\
 双视角语义 + QToken & 80\%\\
 语义 + QToken + QTokenGRU & 87\%\\
 \bottomrule
 \end{tabular}
\end{table}
用户已确认各组使用相同的策略示范数据与训练轮次，
因此该结果是在一致的数据与训练预算下比较不同设计，而不是更多动作示范带来的收益。
这里语义输入不是语义替代全部RGB，具体视觉配置以checkpoint为准。
不能假设各行具有相同分母、训练随机性或相同基础策略，
也不能将约10个百分点的差异直接解释为显著的因果效应。

这些结果提出三项可检验解释。
\emph{表示解释：}语义预训练将对象区域的先验显式提供给策略，
几何监督进一步减少目标关系被编码丢失的风险。
应检查这种收益是否集中于拨空减少、有效接触和完整入区，
并排除单纯增加视角或token容量。
\emph{反馈解释：}当前帧残差能根据最新证据修正当前计划，
因此无需每次重选完整动作块。
应比较相同基础checkpoint下的无残差、MLP残差与频繁ACT查询。
\emph{历史解释：}GRU可借助过去的观测和动作上下文区分当前画面中的运动歧义。
应与历史置乱、仅最新帧及同历史的非循环模型比较，
而不把GRU对MLP的所有收益都归于记忆架构。

从约80\%到87\%还提示几何表示与纠偏可能互补，
但证明完整系统更好不等于证明模块之间存在协同效应。
后续优先以相同平台、视角和预算下的
RGB/语义QToken与无残差/GRU组成匹配的 $2\times2$ 对照，
检验收益是否仍然存在；未完成这一控制前不报告加性分解或协同结论。
理论的职责是把上述解释转为明确条件与反证实验，
不是用已经观察到的收益反推所有假设必然成立。

\subsection{已知任务设置与位置划分}
当前任务为双臂协作暴露目标、将目标完全拨入目标区域、再恢复布料画面占比，
使用三个物体类别和一种布料颜色。
以用户提供的 $(r,c)$ 为网格索引，不将其当作米制坐标，
训练初始位置为
\begin{equation}
 \mathcal P_{\rm tr}=
 \{(0,0),(0,3),(1,1),(1,2),(2,0),(2,3)\}.
\end{equation}
测试覆盖
$\mathcal P_{\rm test}=\{0,1,2\}\times\{0,1,2,3\}$，
其中六个未采集位置为
\begin{equation}
 \mathcal P_{\rm new}=
 \{(0,1),(0,2),(1,0),(1,3),(2,1),(2,2)\}.
\end{equation}
每轮三个物体各测12格，共36个条件单元；
其是否各测一次、各方法实际重复次数和有效样本数尚待日志确认。
应分别报告18个已采集位置条件和18个未采集位置条件的结果，
再报告各物体和总体成功率。

该设计支持同一工作区内未采集初始位置的泛化研究，
不是未见物体类别、布料颜色或任意动力学的泛化验证。
若网格等距且几何布局一致，新位置位于所列四角形成的范围内，
更适合称为工作区内位置泛化，而不是空间外推。
即使初始化未见，示范运动也可能经过这些位置，
因此还不能称为整个状态空间严格未见。
未来若使用另一平台，独立定义该平台的 $\Phi_\ell$，
保留通用输入与纠偏接口，不要求复制相同颜色或网格。

\subsection{从现有结果出发的最小补充验证}
首先固定所有方法的时限、包含判据、覆盖域、覆盖阈值、观察年龄和试验配对，
从日志补齐成功次数与置信区间，按物体和新旧位置分层；
其次在相同基础checkpoint下比较无残差、MLP与GRU，
进行历史置乱、最新观察延迟和动作误差方向诊断，以定位记忆与反馈收益；
随后固定双视角语义基线，对照无query、无几何监督的等量query与有监督query，
同时检查几何精度、动作依赖和拨空/接触/完整入区等事件，
排除视角和容量变化；最后用匹配的 $2\times2$ 表示与GRU对照验证互补性，
仅在主结论稳定后扩展扰动、颜色、数据量或新平台实验。
成功评价采用独立人工复核或独立冻结评估器，
防止使用策略自身分割把同一种错误同时计为正确决策与成功。
所有实验围绕表\ref{tab:preliminary}的机制解释展开，
而不是再同时增加新的阶段头、融合器和控制器。
